\section{Interpretation and unresolved stability questions}

The magnetic calculation gives three distinct types of result.
The fixed-support Hilbert schemes and the isolated dual-ideal
sheaves give numerical DT contributions in a specified
large-volume description. The magnetic extension gives a
one-mode bound-state contribution near its controlled
large-volume wall. The analytically continued periods give
exact phase signs near the quantum transition. Passing from
one statement to another requires stability information;
none is a substitute for the others.

The exact mixed smoothing and selected-base integrability theorems
\cite{Wuebben:2026smoothing} settle existence of the underlying $X_9$
geometric transition. They do not settle deformation of a pair consisting
of the geometry and a magnetic object. Proposition~\ref{prop:probe-smooth-persistence}
does give a local positive result: the two rigid sheaves persist with
their isolated $+1$ contributions in smooth polarized deformations of the
resolved phase, because their obstruction groups vanish. The remaining
problem is passage through the singular transition and physical
stability along the quantum continuation, not an unproved integrability
of the $X_9$ smoothing direction.

In particular the isolated $+1$ contribution at relative
charge $(R,-1)$ does not establish an elementary particle of
charge $R$. The latter class has no connected holomorphic
curve representative in this resolved chamber. Spectral flow
can remove the relative D0 label only while transforming the
core and $B$-field frame as well. The exact phase inequalities
exclude the specified magnetic split and light-particle
emissions near the endpoint, but not every possible decay.
After Higgsing, the nonzero winding of the probe magnetic
charge requires flux strings rather than an isolated particle.

To obtain a physical magnetic spectrum, one would need a
stability condition compatible with the continued integral
periods and control of the other objects at the relevant
charges. Even at large volume, the complete moduli components
have not been classified. The additional reduced-pencil
hypothesis in Appendix~\ref{app:pencil} is another, independent
geometric question. Until these problems are resolved, the
appropriate conclusion is the set of component counts and
wall constraints proved here, not a compact condensate index.
